\begin{trapbox}
	\begin{description}[style=nextline, leftmargin=2.8em, labelsep=0.5em, itemsep=0.35em]
		\item[\textbf{\textcolor{CTrap}{易错点一（非零前提）：}}] 应用 $x_1y_2=x_2y_1$ 时忽略向量非零条件，导致将零向量判定为与任意向量平行。
		\item[\textbf{\textcolor{CTrap}{正确理解（非零限制）：}}] 结论只对非零向量成立。零向量方向任意，不由坐标等式界定。
		\item[\textbf{\textcolor{CTrap}{错因分析（定义不清）：}}] 误以为平行条件不需要非零，实际上坐标推导依赖比例，零向量无法定义比例。
	\end{description}

	\medskip
	\begin{description}[style=nextline, leftmargin=2.8em, labelsep=0.5em, itemsep=0.35em]
		\item[\textbf{\textcolor{CTrap}{易错点二（顺序混淆）：}}] 错误写成 $x_1y_1=x_2y_2$，即对应分量乘积相等。
		\item[\textbf{\textcolor{CTrap}{正确理解（交叉乘积）：}}] 应为 $x_1y_2 = x_2y_1$，即“外乘积”相等，而非内积。
		\item[\textbf{\textcolor{CTrap}{错因分析（记忆混淆）：}}] 与向量垂直条件 $x_1x_2+y_1y_2=0$ 或数量积公式混淆。
	\end{description}

	\medskip
	\begin{description}[style=nextline, leftmargin=2.8em, labelsep=0.5em, itemsep=0.35em]
		\item[\textbf{\textcolor{CTrap}{易错点三（起点不统一）：}}] 判定三点共线时，使用 $\overrightarrow{AB}$ 与 $\overrightarrow{BC}$ 检验，虽可成立但起点不同易导致坐标计算错误或方向混淆。
		\item[\textbf{\textcolor{CTrap}{正确理解（统一起点）：}}] 推荐使用同一起点的向量（如 $\overrightarrow{AB}$ 与 $\overrightarrow{AC}$），避免方向反号等失误。
		\item[\textbf{\textcolor{CTrap}{错因分析（习惯偏差）：}}] 未理解起点统一可简化符号判断，误认为任意两向量平行即能直接判定共线，忽略使用共起点更为稳妥。
	\end{description}
\end{trapbox}
